\chapter{Kiểm chứng và khả năng tái hiện}
\label{chap:verification}
\section{Phân cấp bằng chứng}
Output notebook cung cấp bằng chứng thực nghiệm; JSON/Parquet/NPZ cho phép tái tính; app cung cấp bằng chứng phục vụ; ảnh giao diện xác nhận trạng thái người dùng nhìn thấy. Không thay một loại bằng chứng bằng loại khác: app khởi động được không chứng minh nhãn đúng, và ma trận đúng số lượng không chứng minh corpus không có near-duplicates.
\DataTable{tables/test_summary.tex}{Phạm vi kiểm chứng trên sản phẩm.}

Notebook dữ liệu có 15 cell mã đã thực thi và không có output lỗi. Notebook PhoBERT có 72 cell mã, 65 hình và không có output lỗi. Bốn hình dữ liệu cùng 65 hình mô hình tạo 69 PNG riêng 300 dpi. Mỗi hình trong báo cáo lấy từ output tương ứng, trừ sơ đồ TikZ và ảnh app có nguồn được nêu rõ.

\section{Đối chiếu dữ liệu và kết quả}
Số lượt URL, unique URL, Parquet sạch và corpus sau khử trùng được phân biệt. Số mẫu train + validation + test bằng 199.965. Class weights tính từ train; ma trận có 20 hàng/cột và tổng 29.995; tổng hàng khớp support. Tái tính argmax sau temperature và hệ số trên toàn bộ logits test khớp 29.995 nhãn trong predictions và cùng công thức app.

Mapping id2label/label2id của config, label config, calibration và taxonomy phải đồng nhất. Hash trọng số trong active manifest khớp model. Cache được bảo toàn và không dùng làm dữ liệu runtime của app. Thứ tự mẫu được truy vết qua article index, split manifest và hash văn bản, tránh ghép một nhãn dự đoán với bài khác.

\section{Kiểm tra ứng dụng}
Các kiểm tra bao gồm health, ba vai trò đăng nhập, sai mật khẩu, thiếu phiên, truy cập trái vai trò, nội dung rỗng hoặc chỉ ký hiệu, nhãn ngoài taxonomy, ngưỡng đảo thứ tự, bài không tồn tại, worker Redis--Dramatiq, API của các trang và đăng xuất hủy phiên. Bài minh họa được suy luận thật; không sinh confidence ngẫu nhiên hoặc lấy nhãn tham chiếu thay kết quả model.

Kiểm tra công thức calibration bao gồm nhiệt độ không dương/không hữu hạn và tie theo ID. Frontend được build bằng TypeScript và Vite; giao diện được mở qua trình duyệt để kiểm tra từng vai trò. Chỉ những kiểm tra đã thực thi được ghi là đạt. Không đưa số lượng testcase của một bộ kiểm tra không kèm kết quả vào bằng chứng.

\section{Trình tự tái hiện}
Chạy notebook dữ liệu từ thư mục gốc bằng Python 3.14 để tạo corpus sạch. Chạy notebook PhoBERT với kernel PyTorch CUDA trên GPU NVIDIA. Chế độ fresh bắt đầu fine-tuning từ backbone; các đầu ra được lưu tại dataset, temp, model, result và assets báo cáo. Nội dung HTML có thể thay đổi theo thời điểm, nên crawl lại không cam kết tái tạo đúng số bài của snapshot.

Để xem kết quả mà không train lại, mở output notebook và các file evaluation. Để thực hiện đánh giá lại trong notebook, chọn evaluate với gói tương thích. Để train tiếp, chọn continue và hiểu rằng optimizer/scheduler khởi tạo mới. Những thao tác này khác nhau về mục đích và chi phí; Run All ở chế độ fresh vẫn huấn luyện bảy epoch tối đa.

Ứng dụng khởi động từ gốc bằng Docker Compose; README nêu tài khoản, cổng, backup và log. Báo cáo biên dịch từ source bằng XeLaTeX, BibTeX rồi hai lượt XeLaTeX để giải quyết trích dẫn và tham chiếu. Phông Times New Roman, Cambria Math, khổ A4 và macro trình bày dùng thống nhất theo mẫu.

\section{Giới hạn tái hiện}
Hardware, phiên bản thư viện, precision và kernel có thể thay đổi tốc độ hoặc giá trị số sát ranh giới. Internet ảnh hưởng thu thập và tải backbone. Split seed tái tạo lựa chọn mẫu trong cùng corpus, không thay thế việc khóa snapshot. Metadata môi trường và fingerprint làm rõ phạm vi này; không đưa ra cam kết không có lỗi trong mọi môi trường chưa kiểm tra.
